% No preamble. Prefix fts-. Conditional finite endpoint tuning only.
\section{Finite ultraviolet tuning under an averaged drift hypothesis}
\label{fts-section}

The finite activity induction, leakage estimate, zero-input
parity remainder and real-coupling continuity are available.
This section gives a conditional endpoint argument for their
actual finite local maps. Its additional hypothesis is a
uniformly positive averaged drift of the \emph{actual finite}
leading determinant coefficients. That hypothesis is not
proved here.

\subsection{The actual decrement on the induction envelope}
\label{fts-decrement-section}

Fix the finite-induction data: elementary ratio $R$, grouping
$J$, radius $b$, carrier exponent $a_0$, amplitude block length
$q_a$, endpoint radius $r_b=K_{\mathrm{amp}}b^2$, and intermediate
bound $C_{\mathrm{int}}b^2$. All references and $\beta_j$ are the
actual finite reference-normalized ones. For a real extracted
input define
\begin{align}
 d_j(g,x)&=\beta_{j+1}(F_{j,g}(x)),\notag\\
 \gamma_{j+1}&=\gamma_j-d_j(g_j,x_j),\qquad
 g_j=\gamma_j^{-1/2},\notag\\
 x_{j+1}&=P_{j+1}F_{j,g_j}(x_j).
 \label{fts-actual-update}
\end{align}
Real branches and real initial data give real decrements.
The activity induction applies to every realized weak
coupling sequence while the permitted periods stay above
its threshold. In particular $x_0=0$ is allowed.

Write
\begin{equation}
 b_{0,j}=\beta_{j+1}
       \bigl(D_{M_j,R;N_{j+1}}\bigr)
 \label{fts-leading-coefficient}
\end{equation}
for the finite pivot/Gaussian leading correction. It is
independent of $g$ and of the inherited interaction, but
depends on the actual classical data and regulator. It is
not replaced by an independently selected infinite coefficient.

Since every $x_j$ is extracted, the finite leakage bound,
parity estimate and quantum input comparison give
\begin{align}
 |d_j(g,x_j)-b_{0,j}|
 \leq{}& C_{\mathrm{int}}\omega(N_{\min})
   +C_{\mathrm{par},\beta}
       \{\tau(T(g))+g^2T(g)^6\}\notag\\
 &+C_{\mathrm{in},\beta}\frac{gT(g)}b .
 \label{fts-decrement-error}
\end{align}
The first term is the finite classical leakage, the second
includes the fixed outer-radius factor, and the third includes
$C_{\mathrm{int}}$ and the beta bound. This follows from
$F=T+G$ and $G(x)=G(0)+(G(x)-G(0))$. The $b^2$ in the
inner-radius remainder cancels beta's $b^{-2}$ bound.
No exact finite identity $\beta_{\mathrm{out}}T
=\beta_{\mathrm{in}}$ is used.

At fixed $b,R,J$, any prescribed tolerance $\epsilon>0$
in \eqref{fts-decrement-error} can be obtained by raising
one period threshold and then lowering a coupling ceiling.
The choices are independent of the number of permitted
steps. The finite-induction magnitude bound also supplies
\begin{equation}
                 |d_j(g,x_j)|\leq D_{\mathrm{abs}}
 \label{fts-absolute-increment}
\end{equation}
uniformly on the envelope and in the allowed history lengths
and regulators.

\subsection{The additional hypothesis and positive block increments}
\label{fts-hypothesis-section}

\paragraph{Unproved averaged-drift hypothesis.}
Assume that there are constants $c>0$ and
$C_{\mathrm{fin}}\geq0$, independent of old depth and the
permitted finite periods, such that
\begin{equation}
       \sum_{j=k}^{k+n-1}b_{0,j}\geq cn-C_{\mathrm{fin}}
 \label{fts-drift-hypothesis}
\end{equation}
on the contemplated stage intervals. It suffices to assume
this for the block length and starting stages used below.
No pointwise positivity of an individual coefficient is required.

Choose $q$ to be a positive multiple of $q_a$ with
$C_{\mathrm{fin}}\leq cq/4$. Strengthen the period and
coupling thresholds until the error in
\eqref{fts-decrement-error} is at most $c/4$. Every $q$-step
sequence of projected maps started in the endpoint ball,
at arbitrary prescribed couplings in that weak range, then
satisfies
\begin{equation}
 B_-\leq\sum_{j=k}^{k+q-1}d_j(g_j,x_j)\leq B_+,
 \qquad
 B_-=\frac{cq}{2}>0,\qquad B_+=qD_{\mathrm{abs}}.
 \label{fts-block-decrements}
\end{equation}
Indeed the lower bound is
$cq-C_{\mathrm{fin}}-q\epsilon\geq cq/2$, and the upper
bound follows from \eqref{fts-absolute-increment}.

Each consecutive $q_a$ subblock returns to the same endpoint
ball. Thus the same $C_{\mathrm{int}}b^2$ envelope controls
the entire larger block. This step repeats already proved
quantum blocks; it does not invoke a new direct large-endpoint
high-degree estimate or require a new field radius depending
on $q$.

\subsection{A continuous auxiliary clipped recursion}
\label{fts-clipping-section}

Let $\Gamma_{\mathrm w}=g_{\mathrm w}^{-2}>0$ be the
strengthened weak threshold. For the shooting proof only,
evaluate the map for every real scalar $\gamma$ at
\begin{equation}
 g_{\mathrm{eval}}(\gamma)
       =\max\{\gamma,\Gamma_{\mathrm w}\}^{-1/2},
 \label{fts-clipped-evaluation}
\end{equation}
while leaving the scalar update unclipped:
\begin{align}
 \gamma_{\mathrm{next}}
       &=\gamma-d_j(g_{\mathrm{eval}}(\gamma),x),\notag\\
 x_{\mathrm{next}}
       &=P_{j+1}F_{j,g_{\mathrm{eval}}(\gamma)}(x).
 \label{fts-clipped-update}
\end{align}
This is auxiliary when $\gamma<\Gamma_{\mathrm w}$; no
physical density interpretation is assigned there.
Its evaluated couplings always belong to the proved weak
range. Starting in the endpoint ball, finite induction
keeps all activity arguments admissible independently
of the scalar value.

Joint real-coupling and quotient continuity, together with
continuity of the maximum in
\eqref{fts-clipped-evaluation}, makes every finite clipped
history continuous in its initial $\gamma$. On a compact
initial interval, \eqref{fts-absolute-increment} bounds all
intermediate scalars. Hence their evaluated couplings
belong to a compact positive interval. Neither continuity
at zero coupling nor a continuous representative selection
nor a volume-uniform modulus is needed.

Since \eqref{fts-decrement-error} and
\eqref{fts-drift-hypothesis} apply to every admissible
realized weak sequence, the clipped history also obeys
\eqref{fts-block-decrements}.

\subsection{Endpoint tuning and removal of clipping}
\label{fts-tuning-section}

\begin{theorem}[Conditional finite ultraviolet tuning]
\label{fts-theorem}
Assume \eqref{fts-drift-hypothesis} with the uniformity just
specified. Fix $m\geq1$ and a permitted finite history of
$mq$ steps whose final period is above the strengthened
threshold, with the nested-family compatibility conditions.
Fix a terminal inverse coupling satisfying
\begin{equation}
       \Gamma_*\geq\Gamma_{\mathrm w}+qD_{\mathrm{abs}}.
 \label{fts-terminal-margin}
\end{equation}
For $x_0=0$, or any fixed real extracted input in the
endpoint ball, some initial inverse coupling in
\begin{equation}
       [\,\Gamma_*+mB_-,\ \Gamma_*+mB_+\,]
 \label{fts-shooting-interval}
\end{equation}
generates $mq$ genuine weak local extraction steps and
ends exactly at $\Gamma_*$. The actual scalar relation
is $\gamma_{j+1}=\gamma_j-\beta_{j+1}(F_{j,g_j}(x_j))$,
with $g_j=\gamma_j^{-1/2}$ throughout.
\end{theorem}

\begin{proof}
Run the clipped recursion with the fixed initial activity.
Summing \eqref{fts-block-decrements}, the final scalar
at the left endpoint of \eqref{fts-shooting-interval}
is at most $\Gamma_*$, while at the right endpoint it
is at least $\Gamma_*$. Continuity and the intermediate
value theorem give a clipped history ending at $\Gamma_*$.

For this selected history, sum the positive \emph{future}
block decrements to obtain
\begin{equation}
 \Gamma_*+(m-k)B_-\leq\gamma_{kq}
                   \leq\Gamma_*+(m-k)B_+,
 \qquad 0\leq k\leq m.
 \label{fts-endpoint-bounds}
\end{equation}
Within any block, \eqref{fts-absolute-increment} implies
\begin{equation}
 \gamma_{kq+r}\geq\gamma_{kq}-rD_{\mathrm{abs}}
       \geq\Gamma_*-qD_{\mathrm{abs}}
       \geq\Gamma_{\mathrm w},
 \qquad 0\leq r\leq q.
 \label{fts-no-clipping}
\end{equation}
Thus clipping never occurred on the selected history.
Its scalar and activity updates are the actual local
extraction updates, and its coupling stays weak.
The exact finite representative and vacuum compensation
identities implement the same normalized local density
at every step.
\end{proof}

The initial bare coupling therefore obeys
\begin{equation}
 (\Gamma_*+mB_+)^{-1/2}
       \leq g_{\mathrm{initial}}
       \leq(\Gamma_*+mB_-)^{-1/2}.
 \label{fts-bare-coupling-bounds}
\end{equation}
Conditional on the drift hypothesis, this tunes arbitrarily
long finite histories to one fixed positive weak endpoint,
with bare coupling of order $m^{-1/2}$. It proves no
uniqueness of the tuned initial value or consistency between
the choices at different $m$.

\subsection{What remains unproved}
\label{fts-scope-section}

The additional input is precisely
\eqref{fts-drift-hypothesis} for the actual finite
determinant coefficients with their actual extraction.
The established activity bounds, parity, continuity and
qualitative finite leakage do not prove it. A determinant
cocycle can telescope the infinite counterpart, but a
quantitative one-loop comparison and finite-regulator
control must still yield the required positive averaged
finite drift.

The horizons considered here are multiples of $q$, enough
for an unbounded ultraviolet subsequence if that drift
input is supplied. No assertion about every integer
horizon is needed. The initial period grows to keep the
final period above the common threshold. The theorem
does not establish convergence of the tuned measures,
mixed or large-field histories, a nontrivial continuum
limit, or reconstruction. Its density statement remains
the one for the specified homogeneous local extraction
map and its justified continuation.
